% !TeX program = xelatex
% 来源：../2025年秋季期中试卷B卷答案.pdf（原件未修改）。
% 无答案试卷。独立源文件，使用 XeLaTeX 编译；正文、表格和排版均可编辑。
\documentclass[UTF8,11pt,a4paper,fontset=fandol]{ctexart}
\usepackage[a4paper,top=18mm,bottom=18mm,left=23mm,right=23mm,footskip=9mm]{geometry}
\usepackage{array,tabularx,enumitem,amsmath,amssymb}
\setlength{\parindent}{0pt}
\setlength{\parskip}{0pt}
\linespread{1.12}\selectfont
\pagestyle{plain}
\newcolumntype{C}[1]{>{\centering\arraybackslash}p{#1}}
\newcolumntype{Y}{>{\centering\arraybackslash}X}
\newcommand{\scoretitle}[1]{%
  \noindent\begin{tabular}{|c|c|}\hline
  \raisebox{0pt}[2.4ex][1.4ex]{\bfseries 得分} & \makebox[1.15cm]{} \\ \hline
  \end{tabular}\hspace{1em}\textbf{#1}\par\vspace{0.6em}}
\newcommand{\tfgrid}{%
  \begin{tabularx}{\linewidth}{|*{10}{Y|}}
  \hline 1&2&3&4&5&6&7&8&9&10 \\ \hline
  \rule{0pt}{2.5em} &&&&&&&&& \\ \hline
  \end{tabularx}}
\newcommand{\mcgrid}{%
  \begin{center}\begin{tabularx}{0.74\linewidth}{|*{7}{Y|}}
  \hline 1&2&3&4&5&6&7 \\ \hline
  \rule{0pt}{2.25em} &&&&&& \\ \hline
  \end{tabularx}\end{center}}
\newcommand{\tfitem}[1]{\item #1\nobreak\hfill（\qquad）}
\newcommand{\fourchoices}[4]{%
  \par\nopagebreak\begin{tabularx}{\linewidth}{@{}XXXX@{}}
  A．#1 & B．#2 & C．#3 & D．#4
  \end{tabularx}\par}
\newcommand{\choice}[1]{\par\nopagebreak #1}
\newcommand{\blankpart}[1]{\noindent（#1）\par\vfill}
\begin{document}

\begin{center}
{\small\bfseries 诚实考试吾心不虚，公平竞争方显实力，\\
考试失败尚有机会，考试舞弊前功尽弃。}\\[-0.3em]
\rule{\linewidth}{0.5pt}\\[0.25em]
{\large\bfseries 上海财经大学《宏观经济学》课程期中考试卷（B）\quad 闭卷}\\[0.55em]
课程代码\underline{\makebox[2.6cm]{100096}}\qquad 课程序号\underline{\makebox[2.7cm]{}}\\[0.45em]
{\bfseries 2025---2026 学年第 1 学期}\\[0.7em]
姓名\underline{\makebox[2.3cm]{}}\quad 学号\underline{\makebox[2.6cm]{}}\quad 班级\underline{\makebox[2.3cm]{}}\\[0.65em]
\renewcommand{\arraystretch}{1.2}
\begin{tabular}{|c|*{6}{C{1.5cm}|}}
\hline 题号 & 一 & 二 & 三 & 四 & 五 & 总分 \\ \hline
\rule{0pt}{2.2em}得分 & & & & & & \\ \hline
\end{tabular}
\end{center}
\vspace{-0.3em}
\hspace*{2em}请仔细阅读以下内容，并签名承诺。一旦您签名，将作为诚信承诺的依据，若您在考试过程中出现舞弊行为，将按照学校有关规定严肃处理；若您拒绝进行诚信承诺，将取消考试资格。

\hspace*{2em}考试结束后，考卷必须与答题纸同时上交。未上交考卷，或没有诚信签字的考卷将计以零分。
\begin{center}{\large\bfseries 诚信承诺书}\end{center}
\hspace*{2em}我谨以至诚承诺，在此次考试过程中，我绝无任何舞弊行为，也绝不帮助他人舞弊。

\hfill 学生签名：\underline{\makebox[2.7cm]{}}\par
\hfill \underline{\makebox[1.0cm]{}}年\underline{\makebox[0.6cm]{}}月\underline{\makebox[0.6cm]{}}日\par
\vspace{0.6em}
\scoretitle{一、判断题（每题 2 分，共计 20 分，正确填 $\surd$，错误填 $\times$）}
\tfgrid
\begin{enumerate}[label=\arabic*.,leftmargin=*,itemsep=0.8em,topsep=0.9em]

\tfitem{购买新房既可能是耐用品消费，又可能是投资。}
\tfitem{政府在基础设施建设方面的花费属于投资。}
\tfitem{根据定义，私人储蓄不能为负。}
\end{enumerate}

\begin{enumerate}[label=\arabic*.,start=4,leftmargin=*,itemsep=2.2em,topsep=0.6em]
\tfitem{尽管统计方法有所不同，但 CPI 和 GDP 平减指数是同步且同向变化的。}
\tfitem{人口数量越大的经济体，其人均 GDP 的增长速度就越慢。}
\tfitem{提高失业保险能够缓解失业的痛苦，但同时可能导致结构性失业增加。}
\tfitem{成年人口相同的两个经济体，失业率越低的经济体内就业人数越多。}
\tfitem{持有纸币的机会成本是通货膨胀率。}
\tfitem{由货币乘数产生的货币并没有直接创造任何财富。}
\tfitem{未被预期到的通货膨胀有利于经济中的债务人。}
\end{enumerate}
\newpage
\scoretitle{二、单选题（每小题 4 分，共计 28 分）}
\mcgrid
\begin{enumerate}[label=\arabic*.,leftmargin=*,itemsep=0.55em,topsep=0em]
\item 小王将一万元存入银行账户，该账户的年利率为 5\%。一年后，小王的购买力会（\qquad）
\choice{A．如果通货膨胀率为 $-5\%$，小王的购买力会下降}
\choice{B．如果通货膨胀率为 5\%，小王的购买力不变}
\choice{C．如果通货膨胀率为 10\%，小王的购买力会上升}
\choice{D．以上有多种情形都对}
\item 在一个封闭经济体中，如果当年的消费为 5，投资为 4，政府购买为 3，那么该经济的国民储蓄为（\qquad）
\fourchoices{3}{4}{5}{不确定}
\item 以下哪种商品价格不直接计入 CPI？（\qquad）
\fourchoices{房价}{猪肉价格}{大学学费}{以上都不对}
\item 经济体 A 的实际 GDP 增长率为 4\%，经济体 B 的人均实际 GDP 增长率为 4\%，且两个经济体的人口增长率相同。以下哪项说法正确？（\qquad）
\choice{A．经济体 A 的经济增长速度比经济体 B 快}
\choice{B．经济体 B 的经济增长速度比经济体 A 快}
\choice{C．经济体 A 和 B 的经济增长速度相同}
\choice{D．经济体 A 和 B 的经济增长速度谁快谁慢不确定}
\item 假设税收和政府购买同时增加，且税收增加的幅度小于政府购买增加的幅度。这会导致（\qquad）
\choice{A．均衡下的利率和可贷资金数量均下降}
\choice{B．均衡下的利率和可贷资金数量均上升}
\choice{C．均衡下的利率上升、可贷资金数量下降}
\choice{D．均衡下的利率下降、可贷资金数量上升}
\item 假设经济内找到工作的概率为 10\%，失去工作的概率为 1\%，那么稳态下的失业率为（\qquad）
\fourchoices{1/9}{1/10}{1/11}{以上都不对}
\item 如果央行降低法定存款准备金率，同时在公开市场上买入政府债券，那么货币供给会因此（\qquad）
\fourchoices{增加}{减少}{不变}{不确定}
\end{enumerate}
\newpage
\scoretitle{三、计算题（共 20 分）}
\textbf{1.（本题 20 分）}考虑如下银行 A 的资产负债表
\begin{center}
\renewcommand{\arraystretch}{1.25}
\begin{tabular}{|C{3.8cm}|C{2.1cm}|C{4.8cm}|C{1.5cm}|}
\hline
\multicolumn{2}{|c|}{\textbf{资产}} & \multicolumn{2}{c|}{\textbf{负债和所有者权益}} \\ \hline
存款准备金 & 50 & 存款 & 200 \\ \hline
贷款 & 300 & 一般性债务 & 150 \\ \hline
有价证券 & 50 & 资本（所有者权益） & 50 \\ \hline
\end{tabular}
\end{center}
（1）计算银行 A 的存款准备金率和杠杆率。假设该经济体内的所有银行都选择此存款准备金率，且家庭都不持有任何现金，计算此时的货币乘数。\mbox{（6 分）}

（2）假设该经济体的法定存款准备金率为 15\%，而央行为超额准备金部分支付了一笔 10\% 的利息，计算此情况下该银行所有者权益的变化百分比。\mbox{（7 分）}

（3）假设银行 A 的贷款业务遭受了 20\% 的亏损，\\
\hspace*{2em}计算此情况下该银行所有者权益的变化百分比。\mbox{（7 分）}

\textbf{注：题干内的假设情景互相独立，并非按照先后顺序依次发生。}
\par\medskip\textbf{答：}\par\smallskip
\blankpart{1}\blankpart{2}\blankpart{3}
\newpage
\scoretitle{四、作图分析题（共 20 分）}
\textbf{1.（本题 20 分）}请结合货币供给和需求理论，作图分析

（1）如果央行提高对存款准备金支付的利息，这会对均衡下的货币量和物价水平产生什么影响？请详细解释原因。\mbox{（6 分）}

（2）如果央行提高贴现率，这会对均衡下的货币量和物价水平产生什么影响？\\
\hspace*{2em}请详细解释原因。\mbox{（6 分）}

（3）如果经济内的非流动资产收益率上升，这会对均衡下的货币量和物价水平产生什么影响？请详细解释原因。\mbox{（8 分）}
\par\medskip\textbf{答：}\par\smallskip
\blankpart{1}\blankpart{2}\blankpart{3}
\end{document}
